\documentclass[12pt]{article}
\usepackage{amsmath}
\usepackage{amsfonts}
\usepackage{float}
\usepackage{fancyhdr}
\usepackage{graphicx}
\usepackage{setspace}
\usepackage[colorlinks=true,linkcolor=blue, citecolor=red]{hyperref}
\usepackage{url}
\usepackage[top=.75in, left=.75in, right=.75in, bottom=1in]{geometry}

% For algorithm
\usepackage{algorithm}
\usepackage{algpseudocode}

% ============ CODE ============
\usepackage{listings}
\usepackage{xcolor}
\definecolor{codegreen}{rgb}{0,0.6,0}
\definecolor{codegray}{rgb}{0.5,0.5,0.5}
\definecolor{codepurple}{rgb}{0.58,0,0.82}
\definecolor{backcolour}{rgb}{0.95,0.95,0.92}

% Styling for the code.
\lstdefinestyle{mystyle}{
    backgroundcolor=\color{backcolour},   
    commentstyle=\color{codegreen},
    keywordstyle=\color{magenta},
    numberstyle=\tiny\color{codegray},
    stringstyle=\color{codepurple},
    basicstyle=\ttfamily\footnotesize,
    breakatwhitespace=false,         
    breaklines=true,                 
    captionpos=b,                    
    keepspaces=true,                 
    numbers=left,                    
    numbersep=5pt,                  
    showspaces=false,                
    showstringspaces=false,
    showtabs=false,                  
    tabsize=2
}
\lstset{style=mystyle}

% Table packages and Unicode / Font support
\usepackage{array}
\usepackage{longtable}
\usepackage{enumitem}
\usepackage{booktabs}
\usepackage[utf8]{inputenc}
\DeclareUnicodeCharacter{2264}{\ensuremath{\le}}
\DeclareUnicodeCharacter{2194}{\ensuremath{\leftrightarrow}}
\usepackage[T1]{fontenc}
\usepackage{lmodern}
\graphicspath{{../../../../templates/}{../../prototype/}{figures/}{./}}
\setlength{\emergencystretch}{3em}
\setlist[enumerate]{label=\arabic*.,leftmargin=*}
\setlist[itemize]{leftmargin=*}

% ============ HEADER AND FOOTER ============
\setlength{\headheight}{29.43912pt}

\pagestyle{fancy}
\fancyfoot{}
\fancyfoot[R]{Page \thepage}

\setlength{\parindent}{0pt}
\renewcommand{\baselinestretch}{1.5}

\newcommand{\coursename}{Introduction of Software Engineering}
\newcommand{\reportname}{Studitopia}
\newcommand{\reporttitle}{\href{https://github.com/RoundFloor/Studitopia.git}{Studitopia}}

\newcommand{\studentname}{%
Nguyen Thanh An (24125024)\\
Nguyen Chi Bao (24125025)\\
Tran Nguyen Khanh Binh (24125026)\\
Tran Canh Anh Tuan (24125048)\\
Nguyen Le Quoc Huy (24125100)
}
\newcommand{\teachername}{Ho Tuan Thanh, Nguyen Thi Minh Tuyen}

% Header
\lhead{\reporttitle}
\rhead{
University of Science - VNUHCM\\
\coursename
}

% Footer
\newcommand{\leftfooter}{}
\lfoot{\leftfooter}

% ============ DOCUMENT ============
\begin{document}

\pagenumbering{roman}
\begin{titlepage}
\newcommand{\HRule}{\rule{\linewidth}{0.5mm}}
\centering

\textsc{\LARGE VNU-HCM}\\[0.5cm]
\textsc{\Large University of Science}\\[0.5cm]
\textsc{\large Faculty of Information Technology}\\[0.5cm]
\textsc{Computer Science}\\[0.5cm]

\HRule \\[0.4cm]
{ 
\huge{\bfseries{\reporttitle}}\\[0.5cm]
\large{\bfseries{UC-32 --- Receive Notifications}}
}\\[0.4cm]
\HRule \\[0.5cm]

\textbf{\large Subject: \coursename}\\[0.5cm]

\begin{minipage}[t]{0.55\textwidth}
\begin{flushleft} \large
\emph{Students:}\\
\studentname
\end{flushleft}
\end{minipage}
~
\begin{minipage}[t]{0.4\textwidth}
\begin{flushright} \large
\emph{Supervising Teachers:} \\
\teachername
\end{flushright}
\end{minipage}\\[1cm]

{\large \today}\\[1cm]

\includegraphics[scale=.20]{img/hcmus-logo.png}\\[1cm] 

\vfill
\end{titlepage}	

\tableofcontents
\pagebreak

\pagenumbering{arabic}
\setcounter{page}{1}

\hypertarget{sec:uc-32-receive-notifications}{}
\addcontentsline{toc}{section}{UC-32 --- Receive Notifications}
\section*{UC-32 --- Receive Notifications}
\par\medskip

\textbf{Version:} \texttt{1.0}
\par\medskip

\textbf{Performed by:} [Tran Canh Anh Tuan] | \textbf{Edited by:} [Tran Canh Anh Tuan]
\par\medskip

\begingroup
\small
\renewcommand{\baselinestretch}{1.0}\selectfont
\renewcommand{\arraystretch}{1.15}
\setlength{\tabcolsep}{6pt}
\begin{longtable}{|>{\raggedright\arraybackslash}p{0.09\textwidth}|>{\raggedright\arraybackslash}p{0.84\textwidth}|}
\hline
\textbf{Field} & \textbf{Value} \\ \hline
\endhead
Use Case ID & UC-32 \\ \hline
Name & Receive Notifications \\ \hline
Primary Actor & User \\ \hline
Supporting Actors & Notification Service \\ \hline
Feature Group & Notifications and Safety \\ \hline
Priority & Not explicitly stated \\ \hline
Related User Stories & US-38 \\ \hline
Related Requirements & FR-38.1, FR-38.2, FR-38.3, FR-38.4, FR-38.5, FR-38.6; AC-38.1, AC-38.2, AC-38.3, AC-38.4, AC-38.5, AC-38.6; FR-18.6, FR-20.5, FR-31.4, FR-32.5; NFR-PERF-1, NFR-PERF-6, NFR-SEC-1, NFR-SEC-6, NFR-PRIV-6, NFR-USE-4, NFR-USE-5 \\ \hline
Source Documents & \texttt{proposals/userstory.md} (11.1); \texttt{proposals/functional\_requirement.md} (US-38); \texttt{proposals/study\_room.md}; \texttt{proposals/hall\_room.md}; \texttt{proposals/meeting\_lobby.md}; \texttt{proposals/settings.md}; \texttt{proposals/non\_functional\_requirement.md}; \texttt{proposals/API.md}; \texttt{proposals/response\_code.md}; \texttt{proposals/feature\_flow.md}; \texttt{proposals/screen\_flow.md} \\ \hline
\end{longtable}
\endgroup
\par\medskip

\hypertarget{sec:part-a-use-case-model}{}
\addcontentsline{toc}{subsection}{Part A --- Use Case Model}
\subsection*{Part A --- Use Case Model}
\par\medskip

\hypertarget{sec:a1-focused-use-case-diagram}{}
\addcontentsline{toc}{subsubsection}{A1. Focused Use Case Diagram}
\subsubsection*{A1. Focused Use Case Diagram}
\par\medskip

\begin{center}
\includegraphics[width=0.85\textwidth,height=0.75\textheight,keepaspectratio]{"figures/UC-32-receive-notifications_1.png"}
\end{center}
\par\medskip

\textbf{Diagram notation.} Person shape = human actor; cloud = supporting system; blue = selected use case; solid line = association; labelled solid arrow = generalization; dashed arrow = include or extend.
\par\medskip

\hypertarget{sec:a2-actor-register}{}
\addcontentsline{toc}{subsubsection}{A2. Actor Register}
\subsubsection*{A2. Actor Register}
\par\medskip

\begingroup
\small
\renewcommand{\baselinestretch}{1.0}\selectfont
\renewcommand{\arraystretch}{1.15}
\setlength{\tabcolsep}{6pt}
\begin{longtable}{|>{\raggedright\arraybackslash}p{0.11\textwidth}|>{\raggedright\arraybackslash}p{0.18\textwidth}|>{\raggedright\arraybackslash}p{0.63\textwidth}|}
\hline
\textbf{Actor} & \textbf{Type} & \textbf{Description} \\ \hline
\endhead
User & Primary human actor & A registered user who configures notification types, receives in-app alerts, and marks notifications as read. \\ \hline
Notification Service & Supporting external system actor & Delivers forum, social, help-request, and enabled room-invitation activity to the notification center. \\ \hline
\end{longtable}
\endgroup
\par\medskip

\hypertarget{sec:a3-relationships}{}
\addcontentsline{toc}{subsubsection}{A3. Relationships}
\subsubsection*{A3. Relationships}
\par\medskip

\hypertarget{sec:a3-1-associations}{}
\paragraph*{A3.1 Associations}
\par\medskip

\begingroup
\small
\renewcommand{\baselinestretch}{1.0}\selectfont
\renewcommand{\arraystretch}{1.15}
\setlength{\tabcolsep}{6pt}
\begin{longtable}{|>{\raggedright\arraybackslash}p{0.13\textwidth}|>{\raggedright\arraybackslash}p{0.16\textwidth}|>{\raggedright\arraybackslash}p{0.10\textwidth}|>{\raggedright\arraybackslash}p{0.52\textwidth}|}
\hline
\textbf{Actor} & \textbf{Use Case} & \textbf{Type} & \textbf{Description} \\ \hline
\endhead
User & UC-32 Receive Notifications & Primary & Configures preferences, receives toasts and tray items, and marks notifications read. \\ \hline
Notification Service & UC-32 Receive Notifications & Supporting & Routes enabled activity notifications to the user. \\ \hline
\end{longtable}
\endgroup
\par\medskip

\hypertarget{sec:a3-2-generalization-and-specialization}{}
\paragraph*{A3.2 Generalization and Specialization}
\par\medskip

\begingroup
\small
\renewcommand{\baselinestretch}{1.0}\selectfont
\renewcommand{\arraystretch}{1.15}
\setlength{\tabcolsep}{6pt}
\begin{longtable}{|>{\raggedright\arraybackslash}p{0.09\textwidth}|>{\raggedright\arraybackslash}p{0.09\textwidth}|>{\raggedright\arraybackslash}p{0.09\textwidth}|>{\raggedright\arraybackslash}p{0.65\textwidth}|}
\hline
\textbf{Child} & \textbf{Parent} & \textbf{Type} & \textbf{Inherited Behavior or Association} \\ \hline
\endhead
None & None & None & No actor or use-case generalization is named for US-38 in the proposals. \\ \hline
\end{longtable}
\endgroup
\par\medskip

\hypertarget{sec:a3-3-include}{}
\paragraph*{A3.3 Include}
\par\medskip

\begingroup
\small
\renewcommand{\baselinestretch}{1.0}\selectfont
\renewcommand{\arraystretch}{1.15}
\setlength{\tabcolsep}{6pt}
\begin{longtable}{|>{\raggedright\arraybackslash}p{0.14\textwidth}|>{\raggedright\arraybackslash}p{0.18\textwidth}|>{\raggedright\arraybackslash}p{0.60\textwidth}|}
\hline
\textbf{Base Use Case} & \textbf{Included Use Case} & \textbf{Meaning} \\ \hline
\endhead
None & None & No included use case is named for US-38 in the proposals. \\ \hline
\end{longtable}
\endgroup
\par\medskip

\hypertarget{sec:a3-4-extend}{}
\paragraph*{A3.4 Extend}
\par\medskip

\begingroup
\small
\renewcommand{\baselinestretch}{1.0}\selectfont
\renewcommand{\arraystretch}{1.15}
\setlength{\tabcolsep}{6pt}
\begin{longtable}{|>{\raggedright\arraybackslash}p{0.17\textwidth}|>{\raggedright\arraybackslash}p{0.12\textwidth}|>{\raggedright\arraybackslash}p{0.10\textwidth}|>{\raggedright\arraybackslash}p{0.53\textwidth}|}
\hline
\textbf{Extending Use Case} & \textbf{Base Use Case} & \textbf{Condition} & \textbf{Added Behavior} \\ \hline
\endhead
None & None & None & No extending use case is named for US-38 in the proposals. \\ \hline
\end{longtable}
\endgroup
\par\medskip

\hypertarget{sec:a4-model-notes}{}
\addcontentsline{toc}{subsubsection}{A4. Model Notes}
\subsubsection*{A4. Model Notes}
\par\medskip

\begin{enumerate}
\item Only UC-32 and its source-backed actors are shown; no new I-xx, E-xx, or G-xx identifiers are introduced.
\item The Notification Center is treated as an in-app surface, not as an additional actor; the Notification Service is the supporting external actor.
\item Forum, social, help-request, and room-invitation categories come from \texttt{userstory.md} section 11, \texttt{FR-38.1}--\texttt{FR-38.3}, and \texttt{AC-38.2}--\texttt{AC-38.4}.
\item UC-32 covers receiving and managing notifications; the activities that trigger them remain defined in their own use cases.
\end{enumerate}
\par\medskip

\hypertarget{sec:part-b-use-case-specification}{}
\addcontentsline{toc}{subsection}{Part B --- Use Case Specification}
\subsection*{Part B --- Use Case Specification}
\par\medskip

\hypertarget{sec:b1-brief-description}{}
\addcontentsline{toc}{subsubsection}{B1. Brief Description}
\subsubsection*{B1. Brief Description}
\par\medskip

The user receives real-time in-app toasts and a persistent notification tray with unread indicators for relevant academic and social activity. \texttt{AC-38.1} defines the toast/tray behavior, \texttt{FR-38.4} defines the in-app notification center with unread indicators, and \texttt{AC-38.2}--\texttt{AC-38.4} define the forum, social, and room categories. The user also chooses which notification types trigger alerts and whether audio alerts are enabled, then marks notifications as read (\texttt{AC-38.5}, \texttt{AC-38.6}). \texttt{settings.md} places these controls in \texttt{NotificationSettingsForm} with forum, social, room, and audio controls plus notification-center actions.
\par\medskip

\hypertarget{sec:b2-preconditions}{}
\addcontentsline{toc}{subsubsection}{B2. Preconditions}
\subsubsection*{B2. Preconditions}
\par\medskip

\begin{enumerate}
\item The user is authenticated with an active session (\texttt{NFR-SEC-1}).
\item The user can open \textbf{Settings} and reach the \textbf{Notifications} section.
\item At least one notification type is enabled in the user's notification preferences.
\item A relevant forum, social, help-request, or enabled room-invitation event has occurred.
\item The notification service is available to deliver the event to the in-app center.
\end{enumerate}
\par\medskip

\hypertarget{sec:b3-postconditions}{}
\addcontentsline{toc}{subsubsection}{B3. Postconditions}
\subsubsection*{B3. Postconditions}
\par\medskip

\begin{enumerate}
\item \textbf{Success:} The user's notification preferences are stored with the selected forum, social, room, and audio alert types.
\item \textbf{Success:} An enabled event produces a real-time toast and a persistent notification item with an unread indicator.
\item \textbf{Success:} The user opens the notification center and marks one or more notifications as read.
\item \textbf{Failure:} Disabled event types do not trigger alerts; room invitations are not shown when room notifications are disabled.
\end{enumerate}
\par\medskip

\hypertarget{sec:b4-main-success-scenario-basic-flow}{}
\addcontentsline{toc}{subsubsection}{B4. Main Success Scenario (Basic Flow)}
\subsubsection*{B4. Main Success Scenario (Basic Flow)}
\par\medskip

\begin{enumerate}
\item The user opens \textbf{Settings} and selects \textbf{Notifications}.
\item The system displays \texttt{NotificationSettingsForm} with \texttt{ForumNotificationToggles}, \texttt{SocialNotificationToggles}, \texttt{RoomNotificationToggle}, and \texttt{AudioAlertToggle}.
\item The user chooses which forum, social, and room events trigger alerts and whether audio alerts are enabled.
\item The user selects \textbf{Save Notifications}.
\item The system stores the notification preferences.
\item A relevant enabled event occurs, and the Notification Service routes it to the in-app notification center.
\item The system displays a real-time toast and adds a persistent notification item with an unread indicator.
\item The user selects \textbf{Open Notification Center}.
\item The user reviews the notification tray and selects \textbf{Mark All as Read}, clearing the unread indicators for the marked notifications.
\end{enumerate}
\par\medskip

\hypertarget{sec:b5-alternative-flows}{}
\addcontentsline{toc}{subsubsection}{B5. Alternative Flows}
\subsubsection*{B5. Alternative Flows}
\par\medskip

Unless otherwise stated, an alternative flow resumes at the next step of the basic flow.
\par\medskip

\hypertarget{sec:a1-forum-activity-notification}{}
\paragraph*{A1. Forum activity notification}
\par\medskip

At basic flow step 6, the event is a comment, answer, vote, true-response mark, or mention. The system creates a forum notification under \texttt{FR-38.1} / \texttt{AC-38.2}, provided that the corresponding \texttt{ForumNotificationToggles} choice permits it.
\par\medskip

\hypertarget{sec:a2-social-activity-notification}{}
\paragraph*{A2. Social activity notification}
\par\medskip

At basic flow step 6, the event is a friend request, message, group invitation, or help request. The system creates a social notification under \texttt{FR-38.2} / \texttt{AC-38.3}, provided that the corresponding \texttt{SocialNotificationToggles} choice permits it.
\par\medskip

\hypertarget{sec:a3-room-invitation-notification}{}
\paragraph*{A3. Room invitation notification}
\par\medskip

At basic flow step 6, the event is a room invitation. The system shows it only when \texttt{RoomNotificationToggle} is enabled, as required by \texttt{FR-38.3} / \texttt{AC-38.4}.
\par\medskip

\hypertarget{sec:a4-audio-alert-preference}{}
\paragraph*{A4. Audio alert preference}
\par\medskip

At basic flow step 3, the user enables or disables \texttt{AudioAlertToggle}. The system applies the choice to enabled alerts only and saves it with \textbf{Save Notifications} (\texttt{FR-38.5}, \texttt{AC-38.5}).
\par\medskip

\hypertarget{sec:b6-exception-flows}{}
\addcontentsline{toc}{subsubsection}{B6. Exception Flows}
\subsubsection*{B6. Exception Flows}
\par\medskip

\hypertarget{sec:e1-room-notifications-disabled}{}
\paragraph*{E1. Room notifications disabled}
\par\medskip

A room invitation is generated while \texttt{RoomNotificationToggle} is disabled. The system suppresses the room notification rather than showing it in the tray, toasts, or audio alerts (\texttt{FR-38.3}, \texttt{AC-38.4}).
\par\medskip

\hypertarget{sec:e2-unauthenticated-access-to-notification-settings-or-center}{}
\paragraph*{E2. Unauthenticated access to notification settings or center}
\par\medskip

An unauthenticated request targets Settings or the notification center. The system denies access and requires login (\texttt{NFR-SEC-1}, \texttt{NFR-SEC-6}).
\par\medskip

\hypertarget{sec:b7-business-rules}{}
\addcontentsline{toc}{subsubsection}{B7. Business Rules}
\subsubsection*{B7. Business Rules}
\par\medskip

\begin{enumerate}
\item The system displays real-time toasts and a persistent notification tray with unread indicators (\texttt{AC-38.1}, \texttt{FR-38.4}).
\item Forum notifications include comments, answers, votes, true-response marks, and mentions (\texttt{AC-38.2}, \texttt{FR-38.1}).
\item Social notifications include friend requests, messages, group invitations, and help requests (\texttt{AC-38.3}, \texttt{FR-38.2}).
\item Room invitations are shown only when the user enables room notifications (\texttt{AC-38.4}, \texttt{FR-38.3}).
\item Users choose which notification types trigger alerts and whether audio alerts are enabled (\texttt{AC-38.5}, \texttt{FR-38.5}).
\item Users can mark notifications as read (\texttt{AC-38.6}, \texttt{FR-38.6}).
\end{enumerate}
\par\medskip

\hypertarget{sec:b8-special-requirements}{}
\addcontentsline{toc}{subsubsection}{B8. Special Requirements}
\subsubsection*{B8. Special Requirements}
\par\medskip

\begingroup
\small
\renewcommand{\baselinestretch}{1.0}\selectfont
\renewcommand{\arraystretch}{1.15}
\setlength{\tabcolsep}{6pt}
\begin{longtable}{|>{\raggedright\arraybackslash}p{0.10\textwidth}|>{\raggedright\arraybackslash}p{0.83\textwidth}|}
\hline
\textbf{ID} & \textbf{Requirement} \\ \hline
\endhead
NFR-PERF-6 & Forum notifications must be delivered with a latency of 2 seconds or less under normal operating conditions. \\ \hline
NFR-PERF-1 & Standard user interactions must respond within 2 seconds. \\ \hline
NFR-SEC-1 & The system must require authentication before accessing protected user resources. \\ \hline
NFR-SEC-6 & The system must prevent unauthorized access through server-side authorization checks. \\ \hline
NFR-PRIV-6 & Users must be able to control whether other users can send them messages or room invitations. \\ \hline
NFR-USE-4 & Primary interface actions must use clear labels and recognizable controls. \\ \hline
NFR-USE-5 & The user interface must be usable on screen sizes of at least 320 px in width. \\ \hline
\end{longtable}
\endgroup
\par\medskip

\hypertarget{sec:b9-extension-points}{}
\addcontentsline{toc}{subsubsection}{B9. Extension Points}
\subsubsection*{B9. Extension Points}
\par\medskip

\begingroup
\small
\renewcommand{\baselinestretch}{1.0}\selectfont
\renewcommand{\arraystretch}{1.15}
\setlength{\tabcolsep}{6pt}
\begin{longtable}{|>{\raggedright\arraybackslash}p{0.39\textwidth}|>{\raggedright\arraybackslash}p{0.54\textwidth}|}
\hline
\textbf{Extension Point} & \textbf{Location in Flow of Events} \\ \hline
\endhead
Forum notification preference & Basic flow steps 2-4; alternative A1 \\ \hline
Social notification preference & Basic flow steps 2-4; alternative A2 \\ \hline
Room notification preference & Basic flow steps 2-4; alternatives A3 and E1 \\ \hline
Audio alert preference & Basic flow step 3; alternative A4 \\ \hline
Notification delivery & Basic flow steps 6-7 \\ \hline
Notification-center read actions & Basic flow steps 8-9 \\ \hline
\end{longtable}
\endgroup
\par\medskip

\hypertarget{sec:b10-interface-and-data-notes}{}
\addcontentsline{toc}{subsubsection}{B10. Interface and Data Notes}
\subsubsection*{B10. Interface and Data Notes}
\par\medskip

\begin{itemize}
\item \textbf{Screen/flow path:} \texttt{screen\_flow.md} gives \texttt{Home -$>$ Settings}. \texttt{settings.md} places notification controls under \texttt{SettingsPage} $\rightarrow$ \texttt{NotificationSettingsPage} with suggested route \texttt{/settings/notifications}; the route is a recommendation, not a requirement.
\item \textbf{Notification toggles:} \texttt{ForumNotificationToggles} covers comments, answers, votes, true-response marks, and mentions; \texttt{SocialNotificationToggles} covers friend requests, messages, group invitations, and help requests; \texttt{RoomNotificationToggle} covers optional room invitations; \texttt{AudioAlertToggle} controls sound alerts.
\item \textbf{Notification-center actions:} \texttt{NotificationCenterLink} is labelled \textbf{Open Notification Center}; \texttt{MarkAllReadButton} is labelled \textbf{Mark All as Read}; \texttt{SaveNotificationSettingsButton} is labelled \textbf{Save Notifications}; \texttt{ToastPreview} shows the toast with unread indicator pattern. Individual mark-as-read is derived from \texttt{FR-38.6} / \texttt{AC-38.6}.
\item \textbf{Exact messages:} No quoted notification message text is specified for US-38 in the proposals. The labelled actions above and "Comments, answers, votes, true-response marks, mentions", "Friend requests, messages, group invitations, help requests", and "Optional room invitations" are the source-backed UI descriptions from \texttt{settings.md}.
\item \textbf{Data fields:} \texttt{NotificationPreference} contains \texttt{forumEvents}, \texttt{socialEvents}, \texttt{roomInvitationsEnabled}, and \texttt{audioAlertsEnabled} in \texttt{settings.md}; no notification database schema is specified.
\item \textbf{Notification rules:} Room invitations appear only when room notifications are enabled (\texttt{FR-38.3}, \texttt{AC-38.4}); audio alerts follow \texttt{AC-38.5} / \texttt{FR-38.5}; read state follows \texttt{AC-38.6} / \texttt{FR-38.6}.
\item \textbf{API endpoints/statuses/error codes:} Not specified in proposals. \texttt{API.md} has no notification endpoints, and \texttt{response\_code.md} has no notification error codes; forum codes must not be reused for notifications by analogy.
\end{itemize}
\par\medskip

\hypertarget{sec:b11-source-alignment-notes}{}
\addcontentsline{toc}{subsubsection}{B11. Source Alignment Notes}
\subsubsection*{B11. Source Alignment Notes}
\par\medskip

\begin{itemize}
\item UC-05 step 11 and I-06 Send Forum Notification cover the producer side for comment/answer notifications; UC-32 covers the user's receipt, preferences, and read actions and does not change UC-05's publishing flow.
\item UC-07 and \texttt{study\_room.md} route room invitations and join-request decisions; UC-32 only adds the recipient-side rule that room invitations appear when room notifications are enabled.
\item \texttt{userstory.md} section 11 names the Notification Center and its forum, social, help-request, and room categories; \texttt{settings.md} adds the concrete toggles, labels, and preference fields as UI recommendations derived from \texttt{US-38}.
\item \texttt{API.md} and \texttt{response\_code.md} contain no notification contract; no endpoint, status code, or error code is inferred from the forum APIs.
\item The source does not define retry behavior for failed notification delivery; no retry policy is invented.
\end{itemize}
\par\medskip

\par\medskip
\hypertarget{sec:ui-prototype-specimens}{}
\addcontentsline{toc}{subsection}{UI Prototype Specimens}
\subsection*{UI Prototype Specimens}
\par\medskip
\begin{figure}[H]
\centering
\includegraphics[width=0.75\textwidth,height=0.42\textheight,keepaspectratio]{"../../prototype/settings_04a_settings_notifications_selected.png"}
\par\vspace{2pt}
{\footnotesize \textbf{UI Prototype:} Settings: 04a Settings - Notifications selected}
\end{figure}
\par\smallskip

\end{document}
